% ═══════════════════════════════════════════════════════════════
% LEHRERHINWEISE: Repaso - La ropa y los adjetivos
% Spanisch Klasse 8 - 45 Minuten Einzelstunde
% ═══════════════════════════════════════════════════════════════

\documentclass[11pt, a4paper]{article}

\usepackage[utf8]{inputenc}
\usepackage[T1]{fontenc}
\usepackage[ngerman]{babel}
\usepackage{helvet}
\renewcommand{\familydefault}{\sfdefault}

\usepackage[margin=2cm, top=2cm, bottom=1.5cm]{geometry}
\usepackage{enumitem}
\usepackage{tabularx}
\usepackage{booktabs}
\usepackage{array}
\usepackage{longtable}

\usepackage{xcolor}
\usepackage{tcolorbox}
\tcbuselibrary{skins,breakable}
\usepackage{amssymb}

\usepackage{fancyhdr}
\usepackage{lastpage}

% ═══════════════════════════════════════════════════════════════
% FARBEN
% ═══════════════════════════════════════════════════════════════

\definecolor{spanisch}{HTML}{c41e3a}
\definecolor{grau}{HTML}{666666}
\definecolor{hellgrau}{HTML}{f5f5f5}
\definecolor{loesung}{HTML}{c62828}
\definecolor{basisgreen}{HTML}{2E7D32}
\definecolor{standardblue}{HTML}{1565C0}
\definecolor{erweiterungpurple}{HTML}{7B1FA2}

% ═══════════════════════════════════════════════════════════════
% BEFEHLE
% ═══════════════════════════════════════════════════════════════

\newcommand{\basis}{\textcolor{basisgreen}{\textbf{[B]}}}
\newcommand{\standard}{\textcolor{standardblue}{\textbf{[S]}}}
\newcommand{\erweiterung}{\textcolor{erweiterungpurple}{\textbf{[E]}}}
\newcommand{\lsg}[1]{\textcolor{loesung}{\textbf{#1}}}

\newtcolorbox{hinweisbox}[1][]{
    colback=hellgrau,
    colframe=spanisch,
    boxrule=0.8pt,
    arc=0pt,
    left=8pt, right=8pt, top=6pt, bottom=6pt,
    fonttitle=\bfseries,
    title=#1,
    breakable
}

\newtcolorbox{loesungsbox}[1][]{
    colback=white,
    colframe=loesung,
    boxrule=0.8pt,
    arc=0pt,
    left=8pt, right=8pt, top=6pt, bottom=6pt,
    fonttitle=\bfseries\color{loesung},
    title=#1,
    breakable
}

\setlength{\parindent}{0pt}
\setlength{\parskip}{6pt}

% ═══════════════════════════════════════════════════════════════
% KOPF- UND FUSSZEILE
% ═══════════════════════════════════════════════════════════════

\pagestyle{fancy}
\fancyhf{}
\renewcommand{\headrulewidth}{0.4pt}
\renewcommand{\footrulewidth}{0pt}
\lhead{\small Spanisch Kl. 8}
\chead{\small\textbf{Lehrerhinweise}}
\rhead{\small Seite \thepage/\pageref{LastPage}}

% ═══════════════════════════════════════════════════════════════
% DOKUMENT
% ═══════════════════════════════════════════════════════════════

\begin{document}

% ═══════════════════════════════════════════════════════════════
% TITEL
% ═══════════════════════════════════════════════════════════════

\begin{center}
    {\LARGE\bfseries\textcolor{spanisch}{Lehrerhinweise}}\\[4pt]
    {\large Repaso: La ropa y los adjetivos}\\[2pt]
    {\normalsize 45 Minuten (Einzelstunde) · Klasse 8 · Papierbasiert}
\end{center}

\vspace{8pt}

% ═══════════════════════════════════════════════════════════════
% 1. KURZPLAN
% ═══════════════════════════════════════════════════════════════

\section*{1. Stundenverlauf (45 Minuten)}

\begin{longtable}{|p{1.2cm}|p{2cm}|p{6cm}|p{2.5cm}|p{2cm}|}
\hline
\textbf{Zeit} & \textbf{Phase} & \textbf{Aktivität} & \textbf{Material} & \textbf{SF} \\
\hline
\endhead

0--5 & Einstieg &
Bildimpuls: Person mit Kleidung zeigen. Frage: \textit{¿Qué lleva?} SuS benennen spontan. &
Tafel/Bild &
Plenum \\
\hline

5--13 & Vokabel-revision &
LB austeilen und durchgehen. Kleidung + Adjektive chorisch sprechen. Genus-Regel erklären: -o/-a passt sich an, -e bleibt. &
LB &
Plenum \\
\hline

13--23 & Übung 1 &
AB austeilen. SuS bearbeiten Aufgabe 1+2 (Zuordnung + Lücken). Herumgehen, unterstützen. &
AB &
EA \\
\hline

23--33 & Übung 2 &
Aufgabe 3 (Sätze bilden). Beispiel an Tafel. Differenziert: \basis{} mit Hilfe, \erweiterung{} frei. &
AB, Tafel &
EA \\
\hline

33--40 & Partner-übung &
Partnerphase: ``Beschreibe 2 Kleidungsstücke.'' Redemittel an Tafel. Zeitlimit (5 Min)! &
-- &
PA \\
\hline

40--45 & Sicherung &
2--3 Beispielsätze abfragen. Zusammenfassung: ``Adjektive passen sich an!'' &
Tafel &
Plenum \\
\hline
\end{longtable}

\textbf{Puffer:} Bei Zeitnot Partnerübung auf 3 Min kürzen oder streichen.

% ═══════════════════════════════════════════════════════════════
% 2. DIFFERENZIERUNG
% ═══════════════════════════════════════════════════════════════

\section*{2. Differenzierung}

\begin{tabular}{|l|p{4cm}|p{4cm}|p{4cm}|}
\hline
& \textbf{\basis{} Basis} & \textbf{\standard{} Standard} & \textbf{\erweiterung{} Erweiterung} \\
\hline
\textbf{Aufg. 1} &
Wortliste als Hilfe geben &
Selbstständig &
Zusätzlich: Plural bilden \\
\hline
\textbf{Aufg. 2} &
Erstes Beispiel vorgeben, bei -a/-o helfen &
Selbstständig mit Wortliste &
Ohne Wortliste \\
\hline
\textbf{Aufg. 3} &
Satzanfang vorgeben: \textit{El jersey es...} &
Nach Muster arbeiten &
Freie Sätze ohne Muster \\
\hline
\textbf{Partner} &
Redemittel ablesen erlaubt &
Redemittel als Stütze &
Spontan sprechen \\
\hline
\textbf{Zusatz} &
-- &
Extension bei Zeit &
Extension + eigener Dialog \\
\hline
\end{tabular}

% ═══════════════════════════════════════════════════════════════
% 3. HÄUFIGE FEHLER
% ═══════════════════════════════════════════════════════════════

\section*{3. Häufige Fehler und Reaktionen}

\begin{tabular}{|p{5cm}|p{8.5cm}|}
\hline
\textbf{Fehler} & \textbf{Reaktion} \\
\hline
\textit{*la camiseta bonito} (Endung vergessen) &
Kontrastiv: \textit{el jersey bonit\textbf{o}} vs. \textit{la camiseta bonit\textbf{a}}. Merksatz: ``Der Schatten folgt dem Nomen!'' \\
\hline
\textit{*la falda eleganta} (elegante falsch) &
Regel: Adjektive auf -e ändern sich nicht! \textit{elegante} bleibt \textit{elegante}. \\
\hline
\textit{*los pantalones moderno} (Plural vergessen) &
Auf Plural hinweisen: \textit{los} → \textit{modern\textbf{os}} \\
\hline
\textit{ancho/estrecho} verwechselt &
Gestik: Arme weit = \textit{ancho}, Arme eng = \textit{estrecho} \\
\hline
Unruhe bei Partnerarbeit &
Klares Zeitlimit (Timer), bei Bedarf abbrechen \\
\hline
\end{tabular}

% ═══════════════════════════════════════════════════════════════
% 4. LÖSUNGEN
% ═══════════════════════════════════════════════════════════════

\section*{4. Lösungen AB}

\begin{loesungsbox}[Aufgabe 1: La ropa (6 BE)]
\small
\begin{tabular}{ll@{\hspace{2cm}}ll}
T-Shirt & \lsg{la camiseta} & Hemd & \lsg{la camisa} \\
Pullover & \lsg{el jersey} & Kapuzenpulli & \lsg{la sudadera} \\
Hose & \lsg{los pantalones} & Rock & \lsg{la falda} \\
Kleid & \lsg{el vestido} & Jacke & \lsg{la chaqueta} \\
Mantel & \lsg{el abrigo} & Stiefel & \lsg{las botas} \\
Turnschuhe & \lsg{las zapatillas} & Schuhe & \lsg{los zapatos} \\
\end{tabular}
\end{loesungsbox}

\begin{loesungsbox}[Aufgabe 2: Los adjetivos (8 BE)]
\small
\begin{tabular}{@{}rl@{}}
a) & La camiseta es muy \lsg{cómoda}. \\
b) & El vestido es muy \lsg{elegante}. \\
c) & Los pantalones son muy \lsg{modernos}. \\
d) & La falda es demasiado \lsg{corta}. \\
e) & El jersey es muy \lsg{bonito}. \\
f) & Las zapatillas son muy \lsg{deportivas}. \\
g) & La chaqueta es demasiado \lsg{estrecha}. \\
h) & El abrigo es muy \lsg{largo}. \\
\end{tabular}
\end{loesungsbox}

\begin{loesungsbox}[Aufgabe 3: Sätze bilden (8 BE)]
\small
\begin{tabular}{@{}rl@{}}
a) & \lsg{El jersey es muy bonito.} (oder: Me gusta el jersey porque es bonito.) \\
b) & \lsg{La camisa es muy elegante.} \\
c) & \lsg{Las botas son muy modernas.} \\
d) & \lsg{Los pantalones son demasiado anchos.} \\
\end{tabular}

\textit{Bewertung: 2 BE pro Satz (1 BE Struktur, 1 BE korrekte Adjektiv-Endung)}
\end{loesungsbox}

\begin{loesungsbox}[Aufgabe 4: Mi ropa favorita (3 BE)]
\small
Individuelle Lösung. Beispiel:

\lsg{Mi camiseta favorita es azul. Es muy cómoda y moderna. Llevo mi camiseta favorita porque es bonita.}

\textit{Bewertung: 1 BE pro sinnvollem Satz mit korrekter Adjektiv-Verwendung}
\end{loesungsbox}

\begin{loesungsbox}[Extension Tasks (+5 BE)]
\small
\textbf{E1 (2 BE):}\\
a) ancho $\leftrightarrow$ \lsg{estrecho} \quad b) largo $\leftrightarrow$ \lsg{corto}\\
Beispielsatz: \lsg{La falda es muy corta pero la chaqueta es demasiado larga.}

\vspace{4pt}
\textbf{E2 (3 BE):}\\
Beispieldialog:\\
Cliente: \lsg{Hola, busco un jersey moderno.}\\
Vendedor: \lsg{Aquí tienes. Este jersey es muy bonito.}\\
Cliente: \lsg{Me gusta, pero es demasiado estrecho.}\\
Vendedor: \lsg{Tengo otro más ancho. ¿Te gusta?}
\end{loesungsbox}

% ═══════════════════════════════════════════════════════════════
% 5. MATERIAL-CHECKLISTE
% ═══════════════════════════════════════════════════════════════

\section*{5. Material-Checkliste}

\begin{itemize}[leftmargin=*, itemsep=2pt]
    \item[$\Box$] LB-05-repaso-ropa.pdf (Klassensatz)
    \item[$\Box$] AB-05-repaso-ropa.pdf (Klassensatz)
    \item[$\Box$] Bildimpuls Kleidung (Tafel oder Ausdruck)
    \item[$\Box$] Tafel/Whiteboard für Beispielsätze
    \item[$\Box$] Timer für Partnerphase
\end{itemize}

% ═══════════════════════════════════════════════════════════════
% 6. TAFELANSCHRIEB
% ═══════════════════════════════════════════════════════════════

\section*{6. Tafelanschrieb}

\begin{hinweisbox}[Vorschlag Tafelanschrieb]
\textbf{La ropa y los adjetivos}

\vspace{4pt}
\textbf{Regeln:}
\begin{itemize}[leftmargin=*, itemsep=2pt]
    \item -o/-a: passt sich an! (\textit{el jersey bonit\textbf{o}} / \textit{la camiseta bonit\textbf{a}})
    \item -e: bleibt gleich! (\textit{el vestido elegant\textbf{e}} / \textit{la falda elegant\textbf{e}})
\end{itemize}

\textbf{Redemittel:}
\begin{itemize}[leftmargin=*, itemsep=2pt]
    \item \textit{El/La ... es muy ...}
    \item \textit{Me gusta el/la ... porque es ...}
\end{itemize}
\end{hinweisbox}

% ═══════════════════════════════════════════════════════════════
% 7. HAUSAUFGABE
% ═══════════════════════════════════════════════════════════════

\section*{7. Hausaufgabe (optional)}

\begin{itemize}[leftmargin=*, itemsep=2pt]
    \item Vokabeln wiederholen (Kleidung + Adjektive)
    \item AB fertigstellen (falls nicht geschafft)
    \item \erweiterung{}: Eigenes Outfit beschreiben (3--4 Sätze)
\end{itemize}

\end{document}
